\section{迷宫问题}

\subsection{题型说明}

迷宫问题大致可以分为几类：
\begin{enumerate}

\item \textbf{路径方案总数}:只能使用DFS的回溯法，标记已走路径，回退时撤销标记。
\item \textbf{起点到终点是否可达}：可以用DFS/BFS，只要搜索到终点即可停止。
\item \textbf{最短路径}：只能使用BFS，因为BFS是层序遍历，第一次触碰终点即为最短。
\item \textbf{连通块数量}：可以使用DFS/BFS，最佳方法是遍历地图，每发现一个未标记点就启动一次搜索，并且将本次搜索中遇到的符合条件的点改为不符合条件的点。
	
	
\end{enumerate}

请你思考：
\begin{enumerate}

\item \textbf{什么时候用DFS,什么时候用BFS?}
\item \textbf{什么时候需要用到标记数组,什么时候又不需要用到标记数组？}
\item \textbf{你是否能独立写出DFS和BFS的模板，是否可以根据不同的题型，写出符合条件的代码？}
\item \textbf{是否已经深刻的理解递归的含义，能否掌握方向数组的应用？}


\end{enumerate}

\newpage

\subsection{P1605 迷宫(路径方案总数问题)}

\textbf{没有显式图，需要自己把图给构造出来}

\textbf{题目链接:\href{https://www.luogu.com.cn/problem/P1605}{P1605 迷宫}}

\begin{lstlisting}[language=c++,caption={DFS解法}]
#include<bits/stdc++.h>
using namespace std;
const int N = 10;
int vis[N][N];
int n, m, t;
int sx, sy, fx, fy;
int ax, ay;
int dx[4] = {0, 0, 1, -1};
int dy[4] = {1, -1, 0, 0};
int ans = 0;

void dfs(const int &x, const int &y) {
	if (x == fx && y == fy) {
		ans++;
		return;
	}
	for (int i = 0; i < 4; i++) {
		const int nx = x + dx[i];
		const int ny = y + dy[i];
		if (nx < 1 || nx > n || ny < 1 || ny > m) continue;
		if (!vis[nx][ny]) {
			vis[nx][ny] = true;
			dfs(nx, ny);
			vis[nx][ny] = false;
		}
	}
}

void solve() {
	cin >> n >> m >> t;
	cin >> sx >> sy >> fx >> fy;
	while (t--) {
		cin >> ax >> ay;
		vis[ax][ay] = 1;
	}
	vis[sx][sy] = 1;
	dfs(sx, sy);
	cout << ans << endl;
}

int main() {
	ios::sync_with_stdio(false);
	cin.tie(nullptr);
	solve();
	return 0;
}

\end{lstlisting}

\newpage

\subsection{P1644 跳马问题(路径方案总数问题)}

\textbf{没有显式图，需要自己把图给构造出来，或者用其他方式表示马在图中的状态}

\textbf{题目链接:\href{https://www.luogu.com.cn/problem/P1644}{P1644 跳马问题}}

\begin{lstlisting}[language=c++,caption={DFS解法}]
#include <bits/stdc++.h>
using namespace std;
const int N = 20;
int n, m;
int ans;
int vis[N][N];
int dx[4] = {2, 2, 1, 1};
int dy[4] = {1, -1, 2, -2};

void dfs(const int &x, const int &y) {
	if (x == n && y == m) {
		ans++;
		return;
	}
	for (int i = 0; i < 4; i++) {
		int nx = x + dx[i];
		int ny = y + dy[i];
		if (nx < 0 || nx > n || ny < 0 || ny > m) {
			continue;
		}
		dfs(nx, ny);
	}
}

int main() {
	cin >> m >> n;
	dfs(0, 0);
	cout << ans << endl;
}

\end{lstlisting}

\newpage

\subsection{P1596 水坑计数(连通块数量问题)}

\textbf{重点掌握方向数组的应用}

\textbf{题目链接:\href{https://www.luogu.com.cn/problem/P1596}{P1596 水坑计数}}

\begin{lstlisting}[language=c++,caption={DFS解法以及BFS解法}]
#include<bits/stdc++.h>
using namespace std;
constexpr int N = 101;
char g[N][N];
int n, m, ans = 0;
int dx[8] = {-1, -1, -1, 0, 1, 1, 1, 0};
int dy[8] = {-1, 0, 1, 1, 1, 0, -1, -1};

void dfs(const int &x, const int &y) {
	g[x][y] = '.';
	for (int i = 0; i < 8; i++) {
		int nx = x + dx[i];
		int ny = y + dy[i];
		if (nx < 0 || nx >= n || ny < 0 || ny >= m) {
			continue;
		}
		if (g[nx][ny] == '.') {
			continue;
		}
		dfs(nx, ny);
	}
}

void bfs(const int &x, const int &y) {
	queue<pair<int, int> > q;
	q.emplace(x, y);
	g[x][y] = '.';
	while (!q.empty()) {
		const auto p = q.front();
		q.pop();
		for (int i = 0; i < 8; i++) {
			int nx = p.first + dx[i];
			int ny = p.second + dy[i];
			if (nx < 0 || nx >= n || ny < 0 || ny >= m) {
				continue;
			}
			if (g[nx][ny] == '.') {
				continue;
			}
			g[nx][ny] = '.';
			q.emplace(nx, ny);
		}
	}
}

int main() {
	cin >> n >> m;
	for (int i = 0; i < n; i++) {
		for (int j = 0; j < m; j++) {
			cin >> g[i][j];
		}
	}
	for (int i = 0; i < n; i++) {
		for (int j = 0; j < m; j++) {
			if (g[i][j] == 'W') {
				ans++;
				//dfs(i, j);
				bfs(i, j);
			}
		}
	}
	cout << ans << endl;
}

\end{lstlisting}

\newpage

\subsection{P1451 求细胞数量(连通块数量问题)}

\textbf{每一行的数据如果是连在一起的数字，要用\texttt{scanf}控制每次输入的位数}

\textbf{题目链接:\href{https://www.luogu.com.cn/problem/P1451}{P1451 求细胞数量}}

\begin{lstlisting}[language=c++,caption={DFS解法及BFS解法}]
#include<bits/stdc++.h>
using namespace std;
constexpr int N = 101;
int g[N][N];
int n, m, ans = 0;
int dx[4] = {-1, 0, 1, 0};
int dy[4] = {0, 1, 0, -1};


void dfs(const int &x, const int &y) {
	g[x][y] = 0;
	for (int i = 0; i < 4; i++) {
		int nx = x + dx[i];
		int ny = y + dy[i];
		if (nx < 0 || nx >= n || ny < 0 || ny >= m) {
			continue;
		}
		if (g[nx][ny] == 0) {
			continue;
		}
		dfs(nx, ny);
	}
}

void bfs(const int &x, const int &y) {
	queue<pair<int, int> > q;
	q.emplace(x, y);
	g[x][y] = 0;
	while (!q.empty()) {
		const auto &p = q.front();
		q.pop();
		for (int i = 0; i < 4; i++) {
			int nx = p.first + dx[i];
			int ny = p.second + dy[i];
			if (nx < 0 || nx >= n || ny < 0 || ny >= m) {
				continue;
			}
			if (g[nx][ny] == 0) {
				continue;
			}
			g[nx][ny] = 0;
			q.emplace(nx, ny);
		}
	}
}

int main() {
	cin >> n >> m;
	for (int i = 0; i < n; i++) {
		for (int j = 0; j < m; j++) {
			scanf("%1d", &g[i][j]);
		}
	}
	for (int i = 0; i < n; i++) {
		for (int j = 0; j < m; j++) {
			if (g[i][j] > 0) {
				ans++;
				dfs(i, j);
				//bfs(i, j);
			}
		}
	}
	cout << ans << endl;
}

\end{lstlisting}

\newpage

\subsection{AcWing 844. 走迷宫(最短路径问题)}
\textbf{题目链接:\href{https://www.acwing.com/problem/content/846/}{AcWing 844. 走迷宫}}

\begin{lstlisting}[language=c++,caption={BFS解法}]
#include <bits/stdc++.h>
using namespace std;

const int N = 101;
int n, m;
int g[N][N];
int dist[N][N]; //记录到x,y点的最短路径

int dx[4] = {-1, 0, 1, 0};
int dy[4] = {0, 1, 0, -1};

int bfs(const int &x,const int &y) {
	memset(dist, -1, sizeof(dist)); //将dist全部填充为-1,-1表示没有访问到
	queue<pair<int, int> > q;
	
	q.emplace(x, y);
	dist[x][y] = 0;
	
	while (!q.empty()) {
		const auto &p = q.front();
		q.pop();
		for (int i = 0; i < 4; i++) {
			int nx = p.first + dx[i];
			int ny = p.second + dy[i];
			
			if (nx < 1 || nx > n || ny < 1 || ny > m) continue;
			if (g[nx][ny] == 1) continue;
			if (dist[nx][ny] != -1) continue;
			
			dist[nx][ny] = dist[p.first][p.second] + 1;
			q.emplace(nx, ny);
		}
	}
	return dist[n][m];
}

int main() {
	cin >> n >> m;
	for (int i = 1; i <= n; i++) {
		for (int j = 1; j <= m; j++) {
			cin >> g[i][j];
		}
	}
	
	cout << bfs(1,1) << endl;
	return 0;
}

\end{lstlisting}

\newpage

\subsection{B3625 迷宫寻路(起点到终点是否可达问题)}

\textbf{每一行都是连在一起的字符，可以考虑用\texttt{scanf}直接读入整行的字符串}

\textbf{题目链接:\href{https://www.luogu.com.cn/problem/B3625}{B3625 迷宫寻路}}


\begin{lstlisting}[language=c++,caption={BFS解法}]
#include <bits/stdc++.h>
using namespace std;

constexpr int N = 100;
int n, m;
char g[N][N];
bool vis[N][N];

int dx[4] = {-1, 0, 1, 0};
int dy[4] = {0, 1, 0, -1};

//这个题如果用DFS的话，会超时

void bfs(const int &x, const int &y) {
	queue<pair<int, int> > q;
	q.emplace(x, y);
	vis[0][0] = true;
	
	while (!q.empty()) {
		const auto &p = q.front();
		q.pop();
		
		if (p.first == n - 1 && p.second == m - 1) {
			cout << "Yes" << endl;
			return;
		}
		
		for (int i = 0; i < 4; i++) {
			int nx = p.first + dx[i];
			int ny = p.second + dy[i];
			
			if (nx < 0 || nx >= n || ny < 0 || ny >= m) continue;
			if (vis[nx][ny] || g[nx][ny] == '#') continue;
			
			vis[nx][ny] = true;
			q.emplace(nx, ny);
		}
	}
	
	cout << "No" << endl;
}

int main() {
	cin >> n >> m;
	for (int i = 0; i < n; i++) {
		scanf("%s", g[i]);
	}
	bfs(0, 0);
	return 0;
}

\end{lstlisting}